\section{Risk Analysis}

\subsection{RQ1: Comparative Toxicity Analysis of LLMs}

\begin{itemize}
    \item \textbf{MUST: Quantitative Validation} \textit{RISK:} Low  \textit{IMPORTANCE:} High \\
        We expect that we will be able to qualitatively compare the outputs of the four chosen LLMs and calculate their toxicity scores.
    \item\textbf{SHOULD: } \textit{RISK:} Low  \textit{IMPORTANCE:} High \\
        It is very important to understand why outputs are classified and what the words are that explicitly make the output toxic. We think we will be able to understand and analyse the outputs, as this is the main focus of this project.
    \item \textbf{COULD: Qualitative Validation} \textit{RISK:} High  \textit{IMPORTANCE:} Low \\
        Making a guideline for humans would be very interesting but it is not a priority and falls therefore outside the primary scope of this project. Since this is not very important for the project and it could be quite time-consuming, we do not think we will be able to complete this requirement.
\end{itemize}

\subsection{RQ2: Lexical Analysis of Toxic Inputs}
\begin{itemize}
    \item \textbf{MUST: Toxic Output Subset}
        \textit{RISK:} Low  \textit{IMPORTANCE:} High \\
        It is important to characterise the outputs we consider toxic or not. For this, we have to set a threshold and filter the outputs accordingly to build a dataset of toxic outputs. We think we can complete this requirement as there are already toxicity scoring models that we can use and we only have to set a threshold to filter the outputs.
    
    \item \textbf{MUST: Attribution Explanations} \textit{RISK:} Medium \textit{IMPORTANCE:} High \\
    Highlighting and visualizing what lexical parts make an output toxic is important for understanding how our model works and explaining it. We think it will be a challenge for us to complete this requirement, as we have never used Captum before.

    \item \textbf{MUST: Lexical Analysis}  \textit{RISK:} Medium  \textit{IMPORTANCE:} High \\ By analysing the lexical features of the output, we further understand the model.   Because we do not have a lot of experience with language tools, we think we will have some difficulties completing this requirement.
    
\end{itemize}

\subsection{RQ3: Syntactic Analysis of Toxic Inputs}

\begin{itemize}
    \item \textbf{SHOULD: Syntactic Analysis} \textit{RISK:} Medium  \textit{IMPORTANCE:} Medium  \\
    By analyzing the syntactic features of the output, we further understand the model. Because we do not have a lot of experience with language tools, we think we will have some difficulties completing this requirement.
    \item \textbf{SHOULD: Integration with Lexical Analysis} \textit{RISK:} High \textit{IMPORTANCE:} Low \\
    Tying together the lexical and syntactical analysis creates a better overview of the way the model works, and when and why an output is toxic. However, this is a low priority task, making it less likely that we will complete this requirement.
\end{itemize}